% steps/step14.tex
\section[گام چهاردهم: ملاحظات اخلاقی و قانونی]{\faSync\ \ گام چهاردهم: ملاحظات اخلاقی و قانونی}
\begin{stepbox}

\field{حریم خصوصی کاربران و امنیت داده‌ها}{
    تصاویر پزشکی به شدت حساس هستند و امنیت آن‌ها باید با استانداردهایی نظیر \lr{HIPAA} و \lr{GDPR} مطابقت داشته باشد. فرآیند بی‌نام‌سازی داده‌ها (\lr{De-identification}) و حذف اطلاعات هویتی (\lr{PHI}) از متادیتای فایل‌های \lr{DICOM} باید پیش از ارسال تصویر به وب‌سرویس و در مبدأ انجام شود تا حریم خصوصی بیماران در تمامی مراحل حفظ گردد.
}

\field{شفافیت تصمیم‌گیری و کاهش سوگیری‌ها}{
    جعبه سیاه بودن مدل‌های یادگیری عمیق یک چالش جدی در پزشکی است؛ بنابراین استفاده از ابزارهای هوش مصنوعی قابل‌تفسیر (\lr{XAI}) برای نمایش معیارهای تصمیم‌گیری مدل اهمیت بالایی دارد. علاوه بر این، ارزیابی مدل روی داده‌های توزیع‌های مختلف تضمین می‌کند که سیستم دچار سوگیری (\lr{Bias}) نژادی، جنسیتی یا سنی نشده و عدالت بالینی را رعایت می‌کند.
}

\field{مسئولیت‌پذیری در قبال خطاها و پیامدها}{
    از نظر حقوقی، سیستم‌های هوش مصنوعی تشخیصی در حال حاضر به‌عنوان سیستم‌های پشتیبان تصمیم‌گیری (\lr{Decision Support Systems}) طبقه‌بندی می‌شوند. در سناریوی بروز خطای منفی کاذب (\lr{False Negative}) که به بیمار آسیب جدی می‌رساند، ابزار هوش مصنوعی سلب مسئولیت شده و مسئولیت نهایی تایید یا رد گزارش‌ها و تشخیص نهایی همواره بر عهده پزشک متخصص خواهد بود.
}

\end{stepbox}
